\section{Discussion}
\label{sec:discussion}

\paragraph{A practical rule, with a caveat.}
On CIFAR-100 the results give a simple rule. With an ImageNet-pretrained backbone and up to about 20 images per
class, adding class-prompt images helps, and it helps most at several generated images per real image and on top
of a conventional policy: chosen on validation, such a condition beat the best conventional augmentation in 8 of
10 pretrained cells, by up to 14.9 points. With a randomly initialized classifier, a conventional policy should
come first, and generated images help only on top of it or at a high ratio. With the full training set, generated
images cost accuracy. The caveat is Stage 3. On fine-grained flowers and on dermoscopy images, the recipe that won
on CIFAR-100 lost to conventional augmentation, and LoRA adaptation did not close the gap, although from random initialization it
recovered 5 to 12 points over class-prompt images for ResNet-50 on HAM10000. The rule therefore does not
carry to a new dataset on its own: a short validation pilot against the best conventional policy should decide
whether generated images are worth adding.

\paragraph{What the quality measures miss.}
FID and recall reward a generated set for being close to the real data, and in the scarce-data setting the sets
closest to the budget's real images helped least (Section~\ref{sec:pool-utility}). A set helps by adding what the
few real images lack, provided each added image still shows its class. The two Stage 3 failures break that
proviso in two ways that a distribution-level score summarizes poorly: images that are not of the domain
(HAM10000) and realistic images of the wrong species (Flowers-102). Both are visible at a glance in a per-class
grid of generated images next to real ones, and a per-class measure of class fidelity, such as how well a
classifier trained on the real images recognizes the generated ones (Table~\ref{tab:synthetic-accuracy}), is
closer to what matters than a distance between distributions. We did not test whether such a measure predicts
the gain in Stage 3; it is the measure we would compute first.

\paragraph{What the GAN set adds.}
Stable Diffusion images combine two things: knowledge from outside the dataset, and a pixel distribution that
differs from the training images. The GAN images have neither, because the generator saw only the 5{,}000 real
images, and their effect depends on initialization in the opposite direction: they help classifiers trained from
scratch and hurt pretrained ones. One reading, consistent with the distribution-shift results of
Section~\ref{sec:sign-flip} but not tested directly, is that a pretrained backbone can use images from a shifted
distribution, which add variety the few real images lack, whereas a classifier trained from scratch gains most
from more images like its training images, which the GAN supplies and Stable Diffusion does not. For pretrained
classifiers, the GAN images would then add little that the real images do not already give, while their artefacts
(FID 25.9 against the training images) cost accuracy. The GAN set also shows that the flip is not specific to
Stable Diffusion: whether generated images help depends on how the classifier was initialized for both
generators.

\paragraph{Where the explanations stop.}
Several parts of the explanation are tested. The flip does not come from the diversity or the per-class distance
of the generated images; real-only classifiers trained from scratch recognize the generated images less well; a
classifier trained on generated images alone transfers to real images less well from scratch; and an ensemble
shows that the scratch classifiers still learn something useful from the generated images. Other parts are only
consistent with the data: that the pretrained backbone bridges a distribution shift, that the GAN helps from
scratch because it matches the training images, and that the Stage 3 failures come from class fidelity and the
domain gap rather than from the number of classes, the class imbalance or the metric, which also differ between
the datasets. Each would need its own controlled test, such as filtering generated images by a real-data
classifier's confidence, or a medical dataset with many classes.
